% GENERATED FILE. Edit canonical lessons, then run npm run generate:books.
% canonical-source-hash: fnv1a64:ef53f3f5c39e127a
% canonical-lessons: JA-C513-kikaku-suru, JA-C513-minaosu, JA-C513-kuriageru, JA-C513-chousei-suru, JA-C513-dandori

\chapter{Plans and Projects}
\label{ch:ja-plans-and-projects}
\hlchaptermodality{513}

\begin{chapteropening}
\textbf{By the end of this chapter:} \emph{I can talk about planning and rearranging.}

Complete the last lesson of chapter 513: I can talk about planning and rearranging.
\end{chapteropening}

\section[\texorpdfstring{\ja{きかくする} (kikaku suru)}{きかくする (kikaku suru)}]{\ja{きかくする} (kikaku suru) — to plan (an event); to organize}

\label{lesson:JA-C513-kikaku-suru}



\begin{quote}
Before the new one: say the Japanese for to clash (of plans), then the Japanese for changing every day.
\end{quote}

\subsection*{You'll want to know: \ja{きかくする}}
\textbf{\ja{きかくする}} — \emph{kikaku suru} — \textquotedblleft{}to plan (an event); to organize\textquotedblright{}.

\begin{quote}
\textbf{\ja{もし} \ja{かれが} \ja{りょこうを} \ja{きかくしたら、たのしいでしょう。}} — \emph{moshi kare ga ryokō o kikakushitara, tanoshii deshō} — \textquotedblleft{}If he plans the trip, it will surely be fun.\textquotedblright{}
\end{quote}

Polite \textbf{\ja{きかくします}}, \emph{kikakushimasu}; a \ja{する} verb. Past (polite): \textbf{\ja{かいしゃが} \ja{りょこうを} \ja{きかくしました。}}, \emph{kaisha ga ryokō o kikakushimashita}.

\begin{cousinweb}
\textbf{\ja{きかくする}} (\emph{kikaku suru}) is the Sino-Japanese (kango) noun \emph{kikaku} with the native verb \emph{suru}, \textquotedblleft{}to do\textquotedblright{}: two Chinese characters read with their on'yomi, readings borrowed from Chinese. They are \emph{ki}, \textquotedblleft{}undertake\textquotedblright{} and \emph{kaku}, \textquotedblleft{}brush-stroke\textquotedblright{} (Middle Chinese \emph{hweak}, the reading of this character for \textquotedblleft{}to plan\textquotedblright{}) (KANJIDIC2; Middle Chinese in Baxter's transcription).
\end{cousinweb}

\begin{grammarlens}[title={What you've built}]
One more word for saying what happens if.
\end{grammarlens}

\subsection*{Your turn}
\begin{itemize}
\raggedright
  \item \emph{Say it:} \emph{kikaku suru}
  \item \emph{Say it:} \emph{kikaku suru}, once more
  \item \emph{Say it:} \emph{moshi kare ga ryokō o kikakushitara, tanoshii deshō} — If he plans the trip, it will surely be fun.
  \item \emph{Recall:} say the Japanese for to clash (of plans), then the Japanese for changing every day, then say \emph{kikaku suru} again
\end{itemize}

\begin{tcolorbox}[breakable,colback=teal!4,colframe=teal!35!black,title={Before you move on}]
What does \ja{きかくする} mean? (\textquotedblleft{}To plan (an event); to organize\textquotedblright{}.) Say it once more.
\end{tcolorbox}
\section[\texorpdfstring{\ja{みなおす} (minaosu)}{みなおす (minaosu)}]{\ja{みなおす} (minaosu) — to look over again; to review}

\label{lesson:JA-C513-minaosu}



\begin{quote}
Before the new one: say the Japanese for changing every day, then the Japanese for to plan (an event).
\end{quote}

\subsection*{You'll want to know: \ja{みなおす}}
\textbf{\ja{みなおす}} — \emph{minaosu} — \textquotedblleft{}to look over again; to review\textquotedblright{}.

\begin{quote}
\textbf{\ja{もし} \ja{もんだいが} \ja{あったら、よていを} \ja{みなおします。}} — \emph{moshi mondai ga attara, yotei o minaoshimasu} — \textquotedblleft{}If there is a problem, I will review the schedule.\textquotedblright{}
\end{quote}

Polite \textbf{\ja{みなおします}}, \emph{minaoshimasu}; a godan verb. Past (polite): \textbf{\ja{せんせいが} \ja{しけんを} \ja{みなおしました。}}, \emph{sensei ga shiken o minaoshimashita}.

\begin{cousinweb}
\textbf{\ja{みなおす}} (\emph{minaosu}) is a native Japanese word (wago), inherited rather than borrowed from Chinese, built from \emph{miru}, \textquotedblleft{}to see\textquotedblright{}, and \emph{naosu}, \textquotedblleft{}to fix\textquotedblright{}; its kanji was attached later as a kun'yomi (JMdict, KANJIDIC2).
\end{cousinweb}

\begin{grammarlens}[title={What you've built}]
One more word for saying what happens if.
\end{grammarlens}

\subsection*{Your turn}
\begin{itemize}
\raggedright
  \item \emph{Say it:} \emph{minaosu}
  \item \emph{Say it:} \emph{minaosu}, once more
  \item \emph{Say it:} \emph{moshi mondai ga attara, yotei o minaoshimasu} — If there is a problem, I will review the schedule.
  \item \emph{Recall:} say the Japanese for changing every day, then the Japanese for to plan (an event), then say \emph{minaosu} again
\end{itemize}

\begin{tcolorbox}[breakable,colback=teal!4,colframe=teal!35!black,title={Before you move on}]
What does \ja{みなおす} mean? (\textquotedblleft{}To look over again; to review\textquotedblright{}.) Say it once more.
\end{tcolorbox}
\section[\texorpdfstring{\ja{くりあげる} (kuriageru)}{くりあげる (kuriageru)}]{\ja{くりあげる} (kuriageru) — to bring forward (a date)}

\label{lesson:JA-C513-kuriageru}



\begin{quote}
Before the new one: say the Japanese for to plan (an event), then the Japanese for to look over again.
\end{quote}

\subsection*{You'll want to know: \ja{くりあげる}}
\textbf{\ja{くりあげる}} — \emph{kuriageru} — \textquotedblleft{}to bring forward (a date)\textquotedblright{}.

\begin{quote}
\textbf{\ja{もし} \ja{じかんが} \ja{あったら、かいぎを} \ja{くりあげます。}} — \emph{moshi jikan ga attara, kaigi o kuriagemasu} — \textquotedblleft{}If there is time, we will bring the meeting forward.\textquotedblright{}
\end{quote}

Polite \textbf{\ja{くりあげます}}, \emph{kuriagemasu}; an ichidan verb. Past (polite): \textbf{\ja{がっこうが} \ja{しけんを} \ja{くりあげました。}}, \emph{gakkō ga shiken o kuriagemashita}.

\begin{cousinweb}
\textbf{\ja{くりあげる}} (\emph{kuriageru}) is a native Japanese word (wago), inherited rather than borrowed from Chinese; its kanji was attached later as a kun'yomi, and its older origin is not settled (JMdict, KANJIDIC2).
\end{cousinweb}

\begin{grammarlens}[title={What you've built}]
One more word for saying what happens if.
\end{grammarlens}

\subsection*{Your turn}
\begin{itemize}
\raggedright
  \item \emph{Say it:} \emph{kuriageru}
  \item \emph{Say it:} \emph{kuriageru}, once more
  \item \emph{Say it:} \emph{moshi jikan ga attara, kaigi o kuriagemasu} — If there is time, we will bring the meeting forward.
  \item \emph{Recall:} say the Japanese for to plan (an event), then the Japanese for to look over again, then say \emph{kuriageru} again
\end{itemize}

\begin{tcolorbox}[breakable,colback=teal!4,colframe=teal!35!black,title={Before you move on}]
What does \ja{くりあげる} mean? (\textquotedblleft{}To bring forward (a date)\textquotedblright{}.) Say it once more.
\end{tcolorbox}
\section[\texorpdfstring{\ja{ちょうせいする} (chōsei suru)}{ちょうせいする (chōsei suru)}]{\ja{ちょうせいする} (chōsei suru) — to adjust; to coordinate}

\label{lesson:JA-C513-chousei-suru}



\begin{quote}
Before the new one: say the Japanese for to look over again, then the Japanese for to bring forward (a date).
\end{quote}

\subsection*{You'll want to know: \ja{ちょうせいする}}
\textbf{\ja{ちょうせいする}} — \emph{chōsei suru} — \textquotedblleft{}to adjust; to coordinate\textquotedblright{}.

\begin{quote}
\textbf{\ja{もし} \ja{じかんが} \ja{かわったら、よていを} \ja{ちょうせいします。}} — \emph{moshi jikan ga kawattara, yotei o chōseishimasu} — \textquotedblleft{}If the time changes, I will adjust my schedule.\textquotedblright{}
\end{quote}

Polite \textbf{\ja{ちょうせいします}}, \emph{chōseishimasu}; a \ja{する} verb. Past (polite): \textbf{\ja{せんぱいが} \ja{よていを} \ja{ちょうせいしました。}}, \emph{senpai ga yotei o chōseishimashita}.

\begin{cousinweb}
\textbf{\ja{ちょうせいする}} (\emph{chōsei suru}) is the Sino-Japanese (kango) noun \emph{chōsei} with the native verb \emph{suru}, \textquotedblleft{}to do\textquotedblright{}: two Chinese characters read with their on'yomi, readings borrowed from Chinese. They are \emph{chō}, \textquotedblleft{}harmonize\textquotedblright{} (Middle Chinese \emph{dew}) and \emph{sei}, \textquotedblleft{}organize\textquotedblright{} (\emph{tsyengX}) (KANJIDIC2; Middle Chinese in Baxter's transcription).
\end{cousinweb}

\begin{grammarlens}[title={What you've built}]
One more word for saying what happens if.
\end{grammarlens}

\subsection*{Your turn}
\begin{itemize}
\raggedright
  \item \emph{Say it:} \emph{chōsei suru}
  \item \emph{Say it:} \emph{chōsei suru}, once more
  \item \emph{Say it:} \emph{moshi jikan ga kawattara, yotei o chōseishimasu} — If the time changes, I will adjust my schedule.
  \item \emph{Recall:} say the Japanese for to look over again, then the Japanese for to bring forward (a date), then say \emph{chōsei suru} again
\end{itemize}

\begin{tcolorbox}[breakable,colback=teal!4,colframe=teal!35!black,title={Before you move on}]
What does \ja{ちょうせいする} mean? (\textquotedblleft{}To adjust; to coordinate\textquotedblright{}.) Say it once more.
\end{tcolorbox}
\section[\texorpdfstring{\ja{だんどり} (dandori)}{だんどり (dandori)}]{\ja{だんどり} (dandori) — arrangements; the order of the steps}

\label{lesson:JA-C513-dandori}



\begin{quote}
Before the new one: say the Japanese for to bring forward (a date), then the Japanese for to adjust.
\end{quote}

\subsection*{You'll want to know: \ja{だんどり}}
\textbf{\ja{だんどり}} — \emph{dandori} — \textquotedblleft{}arrangements; the order of the steps\textquotedblright{}.

\begin{quote}
\textbf{\ja{もし} \ja{だんどりが} \ja{よかったら、しごとは} \ja{はやく} \ja{おわるでしょう。}} — \emph{moshi dandori ga yokattara, shigoto wa hayaku owaru deshō} — \textquotedblleft{}If the arrangements are good, the work will probably finish early.\textquotedblright{}
\end{quote}

\begin{cousinweb}
\textbf{\ja{だんどり}} (\emph{dandori}) mixes the two layers, native and Sino-Japanese: Sino-Japanese \emph{dan}, \textquotedblleft{}steps\textquotedblright{} (Middle Chinese \emph{dwanH}), borrowed from Chinese, and native \emph{dori}, \textquotedblleft{}taking\textquotedblright{}, the voiced form of \emph{tori}, from \emph{toru}, \textquotedblleft{}to take\textquotedblright{} (JMdict, KANJIDIC2; Middle Chinese in Baxter's transcription).
\end{cousinweb}

\begin{grammarlens}[title={What you've built}]
One more word for saying what happens if.
\end{grammarlens}

\subsection*{Your turn}
\begin{itemize}
\raggedright
  \item \emph{Say it:} \emph{dandori}
  \item \emph{Say it:} \emph{dandori}, once more
  \item \emph{Say it:} \emph{moshi dandori ga yokattara, shigoto wa hayaku owaru deshō} — If the arrangements are good, the work will probably finish early.
  \item \emph{Recall:} say the Japanese for to bring forward (a date), then the Japanese for to adjust, then say \emph{dandori} again
\end{itemize}

\begin{tcolorbox}[breakable,colback=teal!4,colframe=teal!35!black,title={Before you move on}]
What does \ja{だんどり} mean? (\textquotedblleft{}Arrangements; the order of the steps\textquotedblright{}.) Say it once more.
\end{tcolorbox}
